% Generated draft: requires mathematical and source review.
\documentclass[11pt]{article}
\usepackage[margin=25mm]{geometry}
\usepackage{fontspec}
\setmainfont{Latin Modern Roman}
\usepackage{amsmath,amssymb,mathtools,unicode-math}
\setmathfont{Latin Modern Math}
\usepackage{enumitem,needspace,xurl,hyperref}
\setlength{\parindent}{0pt}
\setlength{\parskip}{5pt}
\setlength{\emergencystretch}{3em}
\newcommand{\N}{\mathbb{N}}
\newcommand{\Z}{\mathbb{Z}}
\newcommand{\Q}{\mathbb{Q}}
\newcommand{\R}{\mathbb{R}}
\newcommand{\C}{\mathbb{C}}
\newcommand{\F}{\mathbb{F}}
\newcommand{\A}{\mathbb{A}}
\newcommand{\PP}{\mathbb{P}}
\newcommand{\E}{\mathbb{E}}
\newcommand{\eps}{\varepsilon}
\newcommand{\dd}{\,\mathrm{d}}
\begin{document}
\section*{UnsolvedMath ID: 30005546}
\textbf{Source problem code:} catalogue entry
\subsection{Definitions and mathematical statement}
For a Cantor set \(X\subset\R\), the supplied record asks for a description of its real-analytic self-maps, including maps into \(X\) that need not be onto, and of its real-analytic automorphism group, denoted in the source by \(\operatorname{Aut}_{\omega}(X)\). Here the usual convention for a Cantor subset of \(\R\) is a nonempty compact perfect nowhere dense set. The supplied excerpt does not define real-analytic maps on such a set or specify the ambient extensions allowed.

To make this definitional issue explicit, one possible convention, for \(X\subset\R^d\), is
\[
\mathcal E_{\omega}(X)=\left\{f:X\to X:\begin{array}{l}\text{there exist an open set }U\subset\R^d\text{ with }X\subset U\\\text{and a real-analytic map }F:U\to\R^d\text{ such that }F|_X=f\end{array}\right\}.
\]
Elements of this set are functions on \(X\), so different extensions having the same restriction represent the same element. A corresponding possible automorphism convention is
\[
\mathcal A_{\omega}(X)=\{f\in\mathcal E_{\omega}(X): f\text{ is bijective and }f^{-1}\in\mathcal E_{\omega}(X)\},
\]
with composition as the group operation. These are illustrative precise conventions, not definitions verified from the report. Requiring extensions to connected neighborhoods, to all of the ambient space, or to ambient analytic diffeomorphisms can change the objects under discussion. The intended definition of \(\operatorname{Aut}_{\omega}(X)\) remains unresolved.

Subject to recovering those definitions, the one-dimensional requests in the supplied record are:
\begin{enumerate}
\item For each Cantor set \(X\subset\R\), characterize its real-analytic self-maps \(f:X\to X\), without assuming surjectivity.
\item Characterize the abstract groups \(G\) for which there exist a Cantor set \(X\subset\R\) and an injective group homomorphism
\[
G\longrightarrow\operatorname{Aut}_{\omega}(X).
\]
This reads ``arise as subgroups'' up to group isomorphism; the original excerpt has damaged wording at this point.
\item Determine which Cantor sets \(X\subset\R\) have finitely generated \(\operatorname{Aut}_{\omega}(X)\); that is, when there exist finitely many elements \(g_1,\ldots,g_n\) such that
\[
\operatorname{Aut}_{\omega}(X)=\langle g_1,\ldots,g_n\rangle.
\]
This request occurs in the supplied clean statement but is absent from the supplied original excerpt.
\end{enumerate}

The higher-dimensional request asks whether an analogue of the report's Theorem 1 holds for \(X\subset\R^d\) and real-analytic maps, under the displayed source condition
\[
\operatorname{int}_X(M\cap X)=\varnothing
\qquad\text{for every real-analytic submanifold }M\subseteq\R^d,
\]
where \(\operatorname{int}_X\) denotes interior in the subspace topology of \(X\). The clean statement interprets the requested conclusion as at most countably many real-analytic maps \(X\to X\).

\textbf{Unresolved qualifications.} The manifold condition, literally read, permits \(M=\R^d\), which makes it impossible for nonempty \(X\). Even a proper open subset of \(\R^d\) is a full-dimensional analytic submanifold, so merely requiring a proper subset does not settle the issue. The intended restriction on submanifolds must be recovered. The map class also needs clarification: every \(x\in X\) gives a distinct constant analytic self-map \(y\mapsto x\), so including all constant maps creates an immediate countability obstruction for a Cantor set. The excerpt does not state Theorem 1, specify whether nonconstant, injective, or surjective maps are intended in its analogue, or explicitly settle whether the higher-dimensional \(X\) must again be a Cantor set. No corrected higher-dimensional conjecture is asserted here.
\subsection{Short English statement}
The record asks which analytic maps can carry a Cantor set into itself, which groups can occur inside its analytic automorphism group, and when that group is finitely generated. It also asks for a higher-dimensional countability analogue. Precise extension conventions and missing restrictions on maps and submanifolds must be recovered before that analogue has a reliable formulation.
\subsection{Sources}
\begin{itemize}
\item Supplied local UnsolvedMath registry record, permanent ID \(30005546\), titled \emph{Analytic Self-Maps and Automorphism Groups of Cantor Sets}; exact source problem code \texttt{catalogue entry}. The record belongs to the Oberwolfach Reports Open Problems collection and assigns the year 2023. No individual proposer is supplied.
\item \emph{Incidence Problems in Harmonic Analysis, Geometric Measure Theory, and Ergodic Theory} (2023), Oberwolfach Reports, pp.~1397--1452, DOI \texttt{10.4171/owr/2023/25}: \url{https://doi.org/10.4171/owr/2023/25}. The supplied original excerpt identifies report page 1425 and refers to ``Problem 3,'' ``Theorem 1,'' and reference [3]; the theorem statement and the bibliographic identity of [3] are not supplied.
\item Supplied literature source, labeled ``OWR 2023 25'': \url{https://publications.mfo.de/bitstream/handle/mfo/4055/OWR_2023_25.pdf?sequence=1}.
\item Literature assessment dated \texttt{2026-08-22}: the supplied curation labels the extensions concerning maps and automorphism groups open and says that the cited theorem proves only a basic countability result. This is supplied dated curation, not an independently verified literature assessment. The record also labels the statement \texttt{corrected\_verified} and reports checking it against the official report; those are supplied verification claims, not verification performed for this draft.
\end{itemize}
\end{document}
